% ---
% Reemplaza líneas 1275-1346 de Docs/Technical_Report.tex
% Decisiones: D-031, H-031
% ---

\subsubsection{\textit{ZZFeatureMap}: evolución de Pauli-Z de segundo orden}

El \textit{ZZFeatureMap} es un circuito de evolución de
Pauli-Z de segundo orden introducido por Havlíček et al.
\cite{havlicek2019supervised} y disponible en Qiskit como
\texttt{zz\_feature\_map}.
Para $n$ qubits y $r$ repeticiones (\texttt{reps}),
el circuito implementa el unitario:

\begin{equation}
  U_{\text{ZZ}}(\mathbf{x}) =
  \left[\mathcal{U}_\Phi(\mathbf{x})\,H^{\otimes n}
  \right]^r
  \label{eq:zzfeaturemap}
\end{equation}

\noindent donde $H^{\otimes n}$ es la capa de compuertas
de Hadamard que inicializa la superposición uniforme, y
$\mathcal{U}_\Phi(\mathbf{x})$ es el bloque de codificación
de datos definido por\footnote{Havlíček et al. y la documentación de Qiskit escriben $+i$; el circuito de Qiskit implementa $-i$, que produce el estado complejo conjugado. El kernel de fidelidad y los valores esperados de este trabajo no dependen de esa elección.}:

\begin{equation}
  \mathcal{U}_\Phi(\mathbf{x}) =
  \exp\!\left(-i \sum_{S \subseteq [n],\,|S|\leq 2}
  \phi_S(\mathbf{x})\,\bigotimes_{j \in S} Z_j\right)
  \label{eq:Uphi}
\end{equation}

\noindent Los coeficientes $\phi_S(\mathbf{x})$ codifican
los datos clásicos en el hamiltoniano de evolución:

\begin{equation}
  \phi_{\{i\}}(\mathbf{x}) = x_i, \qquad
  \phi_{\{i,j\}}(\mathbf{x}) =
  (\pi - x_i)(\pi - x_j)
  \label{eq:phi_coefficients}
\end{equation}

Como $e^{-i\phi Z} = \operatorname{diag}(e^{-i\phi}, e^{i\phi})$, cada término introduce una fase
relativa $2\phi_S$ entre las componentes del estado. Por eso el circuito de Qiskit aplica
compuertas de fase $P(2x_i)$ y $P\bigl(2(\pi-x_i)(\pi-x_j)\bigr)$: el factor 2 corresponde al
parámetro \texttt{alpha} de la función, cuyo valor por defecto es 2. Con $x_i\in[0,\pi]$, la fase
relativa propia recorre la circunferencia una vez.

\noindent El término de un qubit $\phi_{\{i\}}$ codifica
cada característica $x_i$ como una rotación alrededor del
eje $Z$. El término de dos qubits $\phi_{\{i,j\}}$
introduce \textit{entrelazamiento} entre pares de qubits
que codifica las interacciones de segundo orden entre
características, operando mediante la compuerta:

\begin{equation}
  e^{i\,\phi_{\{i,j\}}(\mathbf{x})\,Z_i Z_j} =
  \text{CNOT}_{i \to j}\;
  e^{i\,\phi_{\{i,j\}}(\mathbf{x})\,Z_j}\;
  \text{CNOT}_{i \to j}
  \label{eq:zz_gate}
\end{equation}

\noindent Esta capacidad de capturar interacciones de segundo
orden entre pares de características es la razón por la que
el ZZFeatureMap es uno de los métodos
de codificación de datos más utilizados en tareas de
clasificación cuántica, dada su capacidad para enriquecer la
expresividad del espacio de características cuántico y la
dificultad de simularlo clásicamente. 

Es importante señalar que en Qiskit 2.x el uso de la clase
\texttt{ZZFeatureMap} está marcado como obsoleto
(\textit{deprecated}) a favor de la función
\texttt{zz\_feature\_map()}, que retorna un
\texttt{QuantumCircuit} estándar \cite{ibmquantum2024docs}.
En el presente trabajo se utiliza la función actualizada
compatible con Qiskit 2.x.
